%%% ============================================================
%%% Joint Closure, Per-Coordinate Fuse Data, and a Closed-Form
%%% Algebraic Attractor of Two Commutative Binary Operations on
%%% Z/10Z
%%%
%%% B.R. Sanders, M. Gish (2026)
%%% Target venue: Algebraic Combinatorics
%%% MSC: 20N02, 20N05, 11A07, 11R32
%%%
%%% This is the CONSOLIDATED draft drawing on:
%%%   - the four_core_seed v2 (chain + per-coordinate fuse data)
%%%   - bridge-sprint WP105 (closed-form attractor at alpha=1/2;
%%%       quartic Galois D_4 over LMFDB 4.2.10224.1)
%%%   - bridge-sprint WP110 (4-core fusion-closure structural)
%%%   - bridge-sprint WP113 (alpha-uniqueness PSLQ Stern-Brocot)
%%%
%%% Theorem 1 (printed as Thm thm:chain):    8-element joint-closure chain
%%% Proposition 1 (printed as prop:fuse-data):
%%%                                         per-coordinate fuse data on 4-core
%%% Lemma + Theorems 2-4 (printed as lem:boundary, thm:attractor,
%%%   thm:attractor-half):                  boundary fixed points and
%%%                                         interior existence/uniqueness
%%% Theorem 5 (thm:hbr-ratio):              h/beta = 1+sqrt(3) at alpha=1/2
%%% Theorem 6 (thm:galois):                 R/Br quartic Galois D_4,
%%%                                         LMFDB 4.2.10224.1
%%% Conjecture (conj:alpha-unique):         alpha=1/2 unique rational mixing
%%%                                         weight (PSLQ Stern-Brocot grid)
%%%
%%% Verification: 4core_verification.py + WP105 verification scripts
%%% archived at shared deposit DOI 10.5281/zenodo.18852047.
%%% ============================================================

\documentclass[11pt,reqno]{amsart}

\usepackage{amsmath,amssymb,amsthm,mathtools}
\usepackage[margin=1in]{geometry}
\usepackage[colorlinks=true,linkcolor=blue,citecolor=blue,urlcolor=blue]{hyperref}
\usepackage{microtype}
\usepackage{array}

%% -------- theorem environments --------
\theoremstyle{plain}
\newtheorem{theorem}{Theorem}
\newtheorem{lemma}[theorem]{Lemma}
\newtheorem{proposition}[theorem]{Proposition}
\newtheorem{corollary}[theorem]{Corollary}

\theoremstyle{definition}
\newtheorem{definition}[theorem]{Definition}

\theoremstyle{remark}
\newtheorem*{remark}{Remark}
\newtheorem*{notation}{Notation}

%% -------- macros --------
\newcommand{\Z}{\mathbb{Z}}
\newcommand{\Q}{\mathbb{Q}}
\newcommand{\R}{\mathbb{R}}
\newcommand{\F}{\mathbb{F}}
\newcommand{\Zten}{\Z/10\Z}
\newcommand{\Cfour}{\mathcal{C}}
\newcommand{\Tfuse}{\widehat{T}}
\newcommand{\Bfuse}{\widehat{B}}

\title[Joint closure on $\Zten$]{Joint Closure, Per-Coordinate Fuse
Data, and a Closed-Form Algebraic Attractor of Two Commutative Binary
Operations on $\Zten$}

\author{Brayden R.\ Sanders}
\address{7Site LLC, Hot Springs, AR, USA}
\email{brayden@7site.co}

\author{M.\ Gish}
\address{Independent Researcher, Hot Springs, AR, USA}
\email{monica.gish1992@gmail.com}

\subjclass[2020]{Primary 20N02, 11R32; Secondary 20N05, 11A07, 11Y16}

\keywords{commutative magma, sub-magma chain, joint closure,
fuse-iteration fixed point, quartic number field,
LMFDB 4.2.10224.1, Galois group $D_4$, PSLQ}

\date{\today}

\begin{document}

\begin{abstract}
We study a pair of commutative binary operations
$T,B:\Zten\times\Zten\to\Zten$ given explicitly by their Cayley
tables (both presented in \S\ref{sec:setup}), both arising as restrictions of a common operator-substrate
construction~\cite{SandersGish2026Sigma}. The matrix $T$ used throughout this paper is the upper-triangle authoritative symmetrization of the canonical bit-pattern encoding; $B$ is intrinsically symmetric. Three structural results
emerge from direct enumeration, elementary algebra, and a
finite-dimensional dynamical-systems analysis on the simplex.

(i) The sub-magmas of $\Zten$ that are jointly closed under both
$T$ and $B$ form a strict $8$-element chain in inclusion order,
with sizes $\{1, 4, 5, 6, 7, 8, 9, 10\}$; sizes $2$ and $3$ are
forbidden. Among all $1023$ non-empty subsets of $\Zten$, exactly
these eight satisfy joint closure (Theorem~\ref{thm:chain}). \emph{Lens-scope note: the chain count depends on the symmetrization choice for $T$. The literal bit-pattern (non-symmetrized) variant of $T$ admits a $7$-element chain instead, with size $7$ forbidden, because the asymmetric cells $T[9,4]$ and $T[3,9]$ break closure of the size-$7$ shell $\{0,4,5,6,7,8,9\}$. The $8$-element chain reported in (i) is the canonical symmetrized lens; the 7-element chain on the literal bit-pattern is documented as a parallel lens-dependent result.}

(ii) On the minimal non-trivial element $\Cfour=\{0,7,8,9\}$ of
the chain, the two operations have markedly different
per-coordinate fuse images: $T$ has rank-$2$ image $\{0, 7\}$ on
$\Cfour\times\Cfour$, while $B$ has rank-$4$ image $\Cfour$.
Despite this rank disparity, both fuses preserve $\Cfour$-support
(consequence of the chain), so any convex combination
$F_\alpha = \alpha\Tfuse + (1-\alpha)\Bfuse$ for $\alpha\in[0,1]$
is a self-map of the $4$-core simplex
$\Delta^{3}_{\Cfour}$~(Proposition~\ref{prop:fuse-data}).

(iii) At the symmetric mixing weight $\alpha=\tfrac12$, the
iteration $F_\alpha$ has a unique attractor
$p^*\in\Delta^{3}_{\Cfour}$ with closed-form algebraic ratios:
$p^*_{7}/p^*_{8}=1+\sqrt{3}$ exactly
(Theorem~\ref{thm:hbr-ratio}), and $\xi^*:=p^*_{9}/p^*_{8}$ is
the unique positive real root of the irreducible monic integer
quartic
\[
  f(x) = x^{4}+4x^{3}-x^{2}+2x-2,
\]
whose Galois group over $\Q$ is $D_{4}$ and whose number field
$\Q(\xi^*)$ is LMFDB~$4.2.10224.1$~(Theorem~\ref{thm:galois}).
The polynomial discriminant is
$-2^{6}\cdot 3^{2}\cdot 71$, the field discriminant is
$-2^{4}\cdot 3^{2}\cdot 71$, and $\Q(\sqrt{3})$ is a genuine
subfield of the splitting field, anchoring the $\sqrt{3}$ in
$1+\sqrt{3}$ arithmetically.

A high-precision PSLQ scan over the $17$-point Stern--Brocot grid
of rationals $p/q\in(0,1)$ with $q\le 7$ confirms that
$\alpha=\tfrac12$ is the \emph{unique} rational in this grid at
which both ratios admit small-coefficient algebraic relations
(degree $\le 8$, coefficient $\le 50$);
Conjecture~\ref{conj:alpha-unique} states the corresponding
uniqueness over $\Q\cap(0,1)$ as an open question.

The chain admits a clean structural reading in terms of a
specific permutation $\pi:\Zten\to\Zten$ introduced in
\S\ref{sec:chain}: the chain walks the $\pi$-forward orbit of
$7$, with the $\pi$-fixed elements of $\{0,3,8,9\}$ contributing
at the bottom, at the size-$1\to 4$ jump, and at the size-$7\to 8$
bridge step. (We use $\pi$ rather than $\sigma$ to avoid a name
collision with the non-associativity-rate function $\sigma(N)$
of the companion paper~\cite{SandersGish2026Sigma}.)
\end{abstract}

\maketitle

%% -------- section 1: introduction --------
\section{Introduction}\label{sec:intro}

This paper studies a small object: the family of sub-magmas of
$\Zten$ that are simultaneously closed under two specific
commutative binary operations $T$ and $B$, together with the
fuse-iteration on the resulting $4$-core simplex. We establish
three exact structural results, accompanied by a deterministic
verification script.

The two operations are presented as Cayley tables in
\S\ref{sec:setup}. Both are commutative. Their behavior at the
left-margin is asymmetric: $B$ has $0$ as a two-sided identity
($B(0, j) = j$ for all $j \in \Zten$), while $T$ has $0$ as a
near-zero element with the single off-diagonal exception
$T(0, 7) = 7$ and otherwise $T(0, j) = 0$. Beyond this margin
behavior the operations differ markedly: $T$ has the value $7$
acting as a row absorber on most of $\Zten$; $B$ is approximately
incrementing in the sense that $B(j,j) = j+1 \pmod{10}$ for
$j \in \{1,\dots,7\}$, with diagonal exceptions at
$j \in \{0,8,9\}$ (where $B(0,0)=0$, $B(8,8)=7$, and $B(9,9)=0$).
Despite this surface dissimilarity, three exact structural
phenomena link them tightly: a complete classification of the
joint-closure lattice into a single $8$-element chain
(Theorem~\ref{thm:chain}); the explicit per-coordinate fuse data
on the chain's minimal non-trivial element together with the
rank disparity that distinguishes the two operations
(Proposition~\ref{prop:fuse-data}); and a closed-form algebraic
attractor of the convex-combination iteration
$F_\alpha=\alpha\Tfuse+(1-\alpha)\Bfuse$ on the $4$-core simplex,
exhibited at the symmetric mixing weight $\alpha=\tfrac12$ as
a fixed point in the $D_4$-Galois quartic number field LMFDB~$4.2.10224.1$
(Theorems~\ref{thm:hbr-ratio}~and~\ref{thm:galois}).

\paragraph{The chain (Theorem~\ref{thm:chain}).} The sub-magmas of
$\Zten$ jointly closed under both $T$ and $B$, partially ordered
by inclusion, form a totally ordered $8$-element chain
\[
\{0\}\subset\Cfour\subset\Cfour\cup\{6\}\subset\Cfour\cup\{5,6\}\subset
\Cfour\cup\{4,5,6\}\subset\Cfour\cup\{3,4,5,6\}\subset
\Cfour\cup\{2,3,4,5,6\}\subset\Zten,
\]
where $\Cfour=\{0,7,8,9\}$. The chain is strict: no size-$2$ or
size-$3$ subset of $\Zten$ is jointly closed. Among all $1023$
non-empty subsets of $\Zten$, exactly these eight satisfy joint
closure.

\paragraph{The 4-core fuse data (Proposition~\ref{prop:fuse-data}).}
On the minimal non-trivial element $\Cfour$ of the chain, the two
operations have markedly different per-coordinate fuse images:
$T$ sends mass on $\Cfour\times\Cfour$ into the rank-$2$ image
$\{0, 7\}$, while $B$ sends mass onto the full rank-$4$ image
$\Cfour$. Both fuses preserve $\Cfour$-support (consequence of
joint closure) and both have the same total mass $Z_{T}(p) =
Z_{B}(p) = (p_{0}+p_{7}+p_{8}+p_{9})^{2}$ on $\Cfour$-supported
$p$ --- the latter equality by total-mass conservation: the
bilinear form $Z_{M}(p) = \sum_{i,j}p_{i}p_{j} =
(\sum_{i}p_{i})^{2}$ holds for \emph{any} binary operation $M$.
Together these properties make
$F_\alpha = \alpha\Tfuse + (1-\alpha)\Bfuse$ a continuous family
of self-maps on the $3$-simplex $\Delta^{3}_{\Cfour}$ of
$\Cfour$-supported probability vectors, $\alpha\in[0,1]$.

\paragraph{The closed-form attractor at $\alpha = \tfrac{1}{2}$
(Theorems~\ref{thm:hbr-ratio} and~\ref{thm:galois}).}
The iteration $F_{1/2}$ on $\Delta^{3}_{\Cfour}$ has a unique
attractor $p^*$ with closed-form algebraic ratios. Setting
$(v, h, \beta, r) := (p^*_{0}, p^*_{7}, p^*_{8}, p^*_{9})$, the
fuse-coordinate equations forced by the explicit per-coordinate
form of Proposition~\ref{prop:fuse-data} reduce to
\[
  \left(\frac{h}{\beta}\right)^{2} - 2\,\frac{h}{\beta} - 2 = 0
  \qquad\Longrightarrow\qquad \frac{h}{\beta} = 1 + \sqrt{3},
\]
and the additional constraint linking $r/\beta$ to $h/\beta$
gives that $\xi^* := r/\beta$ is the unique positive real root of
the irreducible monic integer quartic
\[
  f(x) = x^{4} + 4x^{3} - x^{2} + 2x - 2,
\]
whose Galois group over $\Q$ is $D_{4}$. The number field
$K = \Q[x]/(f)$ has signature $(2, 1)$, class number $1$,
discriminant $d_{K} = -2^{4}\cdot 3^{2}\cdot 71$, and is the
catalogued field LMFDB~$4.2.10224.1$. The splitting field of $f$
over $\Q$ has degree $8$, signature $(0, 4)$, class number $14$,
and discriminant $2^{8}\cdot 3^{4}\cdot 71^{4}$. The intermediate
field $\Q(\sqrt{3})$ is a genuine subfield, anchoring the
$\sqrt{3}$ in the closed-form $h/\beta = 1+\sqrt{3}$:
\[
  f(x) = \bigl(x^{2}+(2-\sqrt{3})x+(\sqrt{3}-1)\bigr)
         \bigl(x^{2}+(2+\sqrt{3})x-(\sqrt{3}+1)\bigr).
\]

\paragraph{$\alpha$-uniqueness on the Stern--Brocot grid
(Conjecture~\ref{conj:alpha-unique}).} A high-precision PSLQ
scan~\cite{FergusonBailey1999} over the $17$-point Stern--Brocot
grid of rationals $p/q\in(0,1)$ with $q\le 7$ confirms that
$\alpha=\tfrac12$ is the only rational in this grid at which both
$h/\beta$ and $r/\beta$ admit small-coefficient algebraic
relations (degree $\le 8$, integer coefficients $\le 50$ in
absolute value). The corresponding uniqueness over
$\Q\cap(0,1)$ is stated as Conjecture~\ref{conj:alpha-unique}.

\paragraph{The pair $(T,B)$ is not independent.} Both Cayley
tables arise from a common operator-substrate construction in
which $T$ is the rank-3 threshold projection at threshold
$T^{*} = 5/7$ and $B$ is the full rank-10 composition; the
construction is detailed in the companion
paper~\cite{SandersGish2026Sigma}. The present paper takes
$(T, B)$ at $N = 10$ as given and studies the joint-closure
structure of the pairing. The chain rigidity above has been
verified for analogous tables $(T_{p}, B_{p})$ constructed via
the same operator-substrate recipe over $\F_{p}$-substrates of
comparable size, for $p \in \{2, 3, 5, 7, 11, 13\}$; the
construction of $(T_{p}, B_{p})$ from $\F_{p}$ and the resulting
$\F_{p}$-universality of the joint-closure structure are the
subject of a companion paper in
preparation~\cite{SandersClaudeChat2026BridgeSprint}.

\paragraph{Positioning.} The classical Drápal--Kepka tradition
of low-density associativity in
quasigroups~\cite{Kepka1980, DrapalKepka1985, DrapalWanless2021}
treats non-associativity counts on cancellative structures. The
present paper studies the joint closure of two non-cancellative
magmas, the per-coordinate fuse data on the resulting $4$-core,
and a closed-form algebraic attractor of the convex-combination
iteration on the $4$-core simplex. The simplex-iteration
component of the present results sits within the broader theory
of positive bilinear maps and contraction mappings in projective
metrics~\cite{Birkhoff1957, Bushell1973}; we use the explicit
algebraic structure of $(T, B)$ to extract closed-form
fixed-point coordinates rather than invoke that theory's
asymptotic-rate machinery directly. The closed-form attractor
identifies an explicit integer quartic whose Galois closure
recovers the catalogued number field LMFDB~$4.2.10224.1$ as the
ring of definition of a fuse-iteration fixed point in
$\Delta^{3}_{\Cfour}$; to our knowledge this connection has not
appeared previously in the finite-magma or quasigroup literature.

\paragraph{Outline.} \S\ref{sec:setup} fixes the two tables.
\S\ref{sec:chain} proves the joint-closure chain
(Theorem~\ref{thm:chain}). \S\ref{sec:4core} records the
corollary that $\Cfour=\{0,7,8,9\}$ is the minimal non-trivial
element of the chain. \S\ref{sec:norm} records the per-coordinate
fuse data on $\Cfour$ (Proposition~\ref{prop:fuse-data}) with
the rank disparity between $\Tfuse$ and $\Bfuse$.
\S\ref{sec:dynamics} introduces the convex-combination iteration
$F_\alpha$ on $\Delta^{3}_{\Cfour}$ and proves global
attraction of a unique fixed point for $\alpha\in(0,1)$.
\S\ref{sec:hbr} derives the closed-form ratio $h/\beta = 1+\sqrt 3$
at $\alpha=\tfrac12$ (Theorem~\ref{thm:hbr-ratio}).
\S\ref{sec:galois} derives the integer quartic $f$ for $\xi^* =
r/\beta$, computes its Galois group as $D_{4}$, and identifies
the number field $K = \Q(\xi^*)$ as LMFDB~$4.2.10224.1$
(Theorem~\ref{thm:galois}). \S\ref{sec:alpha-unique} reports the
PSLQ Stern--Brocot scan and states the
$\alpha$-uniqueness conjecture
(Conjecture~\ref{conj:alpha-unique}).
\S\ref{sec:verify} catalogs the verification scripts.
\S\ref{sec:scope} states the limitations of the present results.
\S\ref{sec:forward} lists companion papers in preparation
that develop the dynamical and number-theoretic structure of the
pair $(T,B)$.

%% -------- section 2: setup --------
\section{Definitions and the two tables}\label{sec:setup}

We work throughout on the additive group $\Zten$, identified with
the labels $\{0,1,\dots,9\}$.

\begin{definition}[The two tables]\label{def:tables}
The two binary operations $T,B:\Zten\times\Zten\to\Zten$ are given
by their Cayley tables:
\[
T = \left[
\begin{array}{*{10}{c}}
0 & 0 & 0 & 0 & 0 & 0 & 0 & 7 & 0 & 0 \\
0 & 7 & 3 & 7 & 7 & 7 & 7 & 7 & 7 & 7 \\
0 & 3 & 7 & 7 & 4 & 7 & 7 & 7 & 7 & 9 \\
0 & 7 & 7 & 7 & 7 & 7 & 7 & 7 & 7 & 3 \\
0 & 7 & 4 & 7 & 7 & 7 & 7 & 7 & 8 & 7 \\
0 & 7 & 7 & 7 & 7 & 7 & 7 & 7 & 7 & 7 \\
0 & 7 & 7 & 7 & 7 & 7 & 7 & 7 & 7 & 7 \\
7 & 7 & 7 & 7 & 7 & 7 & 7 & 7 & 7 & 7 \\
0 & 7 & 7 & 7 & 8 & 7 & 7 & 7 & 7 & 7 \\
0 & 7 & 9 & 3 & 7 & 7 & 7 & 7 & 7 & 7
\end{array}
\right],\qquad
B = \left[
\begin{array}{*{10}{c}}
0 & 1 & 2 & 3 & 4 & 5 & 6 & 7 & 8 & 9 \\
1 & 2 & 3 & 4 & 5 & 6 & 7 & 2 & 6 & 6 \\
2 & 3 & 3 & 4 & 5 & 6 & 7 & 3 & 6 & 6 \\
3 & 4 & 4 & 4 & 5 & 6 & 7 & 4 & 6 & 6 \\
4 & 5 & 5 & 5 & 5 & 6 & 7 & 5 & 7 & 7 \\
5 & 6 & 6 & 6 & 6 & 6 & 7 & 6 & 7 & 7 \\
6 & 7 & 7 & 7 & 7 & 7 & 7 & 7 & 7 & 7 \\
7 & 2 & 3 & 4 & 5 & 6 & 7 & 8 & 9 & 0 \\
8 & 6 & 6 & 6 & 7 & 7 & 7 & 9 & 7 & 8 \\
9 & 6 & 6 & 6 & 7 & 7 & 7 & 0 & 8 & 0
\end{array}
\right].
\]
Both $T$ and $B$ are commutative.
\end{definition}

\begin{notation}
We refer to the indices $0,7,8,9$ frequently. In proofs we
sometimes write $(p_{0},p_{7},p_{8},p_{9})=(v,h,\beta,r)$ for
masses on these four indices.
\end{notation}

\begin{definition}[Sub-magma; joint closure]\label{def:submagma}
A subset $S\subseteq\Zten$ is a \emph{sub-magma} of $(\Zten,M)$ if
$M(i,j)\in S$ for all $i,j\in S$. We call $S$ \emph{jointly
closed} (under $T$ and $B$) if it is a sub-magma of both. The set
of jointly closed subsets, ordered by inclusion, is denoted
$\mathcal{L}(T,B)$.
\end{definition}

\begin{definition}[Fuse and normalizer]\label{def:fuse}
For a binary operation $M:\Zten\times\Zten\to\Zten$ and a vector
$p=(p_{0},\dots,p_{9})\in\R^{10}$, the \emph{fuse}
$\widehat{M}(p)\in\R^{10}$ is the convolution along the operation:
\[
\widehat{M}(p)_{c}\;=\;\sum_{(i,j)\in\Zten^{2}\,:\,M(i,j)=c}p_{i}p_{j},
\qquad c\in\Zten.
\]
The \emph{normalizer} is $Z_{M}(p)=\sum_{c\in\Zten}\widehat{M}(p)_{c}$.
\end{definition}

\begin{remark}\label{rem:total-mass}
By bilinearity, $Z_{M}(p) = (\sum_{i}p_{i})^{2}$ identically for
any binary operation $M:\Zten\times\Zten\to\Zten$ and any
$p\in\R^{10}$, since each ordered pair $(i,j)\in\Zten^{2}$
contributes $p_{i}p_{j}$ to exactly one output bucket. In
particular, for $p$ supported on $\Cfour$ we have
$\sum_{i}p_{i} = v+h+\beta+r$, and so
$Z_{T}(p) = Z_{B}(p) = (v+h+\beta+r)^{2}$ holds by elementary
total-mass conservation, with no input from the joint-closure
structure. We state this explicitly and direct the reader to
Proposition~\ref{prop:fuse-data} below for the substantive
content of the $\Cfour$-fuse data, which is the
\emph{per-coordinate} symbolic form of $\Tfuse(p)$ and
$\Bfuse(p)$ and the structural rank disparity between them.
\end{remark}

%% -------- section 3: chain --------
\section{The joint-closure chain}\label{sec:chain}

\begin{theorem}[Joint-closure chain]\label{thm:chain}
The sub-magmas of $\Zten$ jointly closed under both $T$ and $B$
form a strict chain of length $8$ in inclusion order:
\[
\mathcal{L}(T,B)\;=\;\bigl\{S_{1}\subset S_{4}\subset S_{5}\subset
S_{6}\subset S_{7}\subset S_{8}\subset S_{9}\subset S_{10}\bigr\},
\]
where
\begin{align*}
S_{1}&=\{0\}, & S_{4}&=\{0,7,8,9\}, \\
S_{5}&=\{0,6,7,8,9\}, & S_{6}&=\{0,5,6,7,8,9\}, \\
S_{7}&=\{0,4,5,6,7,8,9\}, & S_{8}&=\{0,3,4,5,6,7,8,9\}, \\
S_{9}&=\{0,2,3,4,5,6,7,8,9\}, & S_{10}&=\Zten.
\end{align*}
The chain is strict: no subset of $\Zten$ of size $2$ or size $3$
is jointly closed.
\end{theorem}

\begin{proof}
We verify the joint-closure status of each candidate subset by
direct computation. The proof has two parts: (a) confirming that
each $S_{k}$ in the displayed list is jointly closed, and (b)
confirming that no other non-empty subset is jointly closed.

\textbf{Part (a) --- the eight subsets are jointly closed.} For
each $k\in\{1,4,5,6,7,8,9,10\}$, the verification reduces to
checking that $T(i,j)\in S_{k}$ and $B(i,j)\in S_{k}$ for all
$i,j\in S_{k}$. We display the size-$4$ case explicitly and refer
to the verification script (\S\ref{sec:verify}) for the others.

For $S_{4}=\{0,7,8,9\}$, the $4\times 4$ restrictions are
\[
T\big|_{S_{4}\times S_{4}}=\begin{pmatrix}
0 & 7 & 0 & 0 \\
7 & 7 & 7 & 7 \\
0 & 7 & 7 & 7 \\
0 & 7 & 7 & 7
\end{pmatrix},
\qquad
B\big|_{S_{4}\times S_{4}}=\begin{pmatrix}
0 & 7 & 8 & 9 \\
7 & 8 & 9 & 0 \\
8 & 9 & 7 & 8 \\
9 & 0 & 8 & 0
\end{pmatrix}.
\]
Every entry of $T\big|_{S_{4}\times S_{4}}$ lies in
$\{0,7\}\subseteq S_{4}$. Every entry of
$B\big|_{S_{4}\times S_{4}}$ lies in $\{0,7,8,9\}=S_{4}$.

\textbf{Part (b) --- no other subset is jointly closed.}

\emph{Sizes $2$ and $3$.} We dispatch these by exhaustive case
analysis on the diagonal entries of $B$. Inspecting the diagonal:
\[
B(j,j)=\begin{cases}
0 & j=0,\\
j+1 & j\in\{1,2,3,4,5,6,7\},\\
7 & j=8,\\
0 & j=9.
\end{cases}
\]

\emph{Size $2$.} Let $S=\{a,b\}$ with $a<b$. Joint closure requires
$B(a,a),B(b,b)\in S$ and $T(a,a),T(b,b)\in S$, plus the off-diagonal
constraints.

\emph{Subcase $0\in S$.} Then $S=\{0,a\}$ with $a\in\{1,\dots,9\}$.
The diagonal $B(a,a)$ must lie in $\{0,a\}$. From the diagonal
table: for $a\in\{1,\dots,7\}$, $B(a,a)=a+1\notin\{0,a\}$; for
$a=8$, $B(8,8)=7\notin\{0,8\}$; for $a=9$, $B(9,9)=0\in\{0,9\}$ and
$B(0,9)=B(9,0)=9\in\{0,9\}$, so $\{0,9\}$ is $B$-closed. However
$T(9,9)=7\notin\{0,9\}$, so $\{0,9\}$ fails $T$-closure.

\emph{Subcase $0\notin S$.} Then $S\subseteq\{1,\dots,9\}$, and
both $B(a,a)\in S$ and $B(b,b)\in S$ are required. Since
$B(j,j)\neq j$ for all $j\in\{1,\dots,9\}$, we need $B(a,a)=b$ and
$B(b,b)=a$. Listing all $j\in\{1,\dots,9\}$ with their diagonal
values: $(1,2),(2,3),(3,4),(4,5),(5,6),(6,7),(7,8),(8,7),(9,0)$.
The only unordered pair $\{a,b\}\subseteq\{1,\dots,9\}$ with
$B(a,a)=b$ and $B(b,b)=a$ is $\{7,8\}$ (since $B(7,7)=8$ and
$B(8,8)=7$). However $B(7,8)=9\notin\{7,8\}$, so $\{7,8\}$ fails
$B$-closure.

The remaining size-$2$ candidates do not satisfy
$B(a,a),B(b,b)\in S$ and are excluded immediately. No size-$2$
subset is jointly closed.

\emph{Size $3$.} Suppose $S=\{a,b,c\}$ with $a<b<c$ is jointly
closed. The three diagonal entries $B(a,a),B(b,b),B(c,c)$ and the
three diagonals $T(s,s)$ for $s\in S$ must all lie in $S$. Direct
enumeration of $\binom{10}{3}=120$ candidates against just the
six diagonal constraints leaves only two candidates:
$\{0,7,8\}$ and $\{6,7,8\}$. (For $\{0,7,8\}$: $B$-diagonals are
$0,8,7$; for $\{6,7,8\}$: $B$-diagonals are $7,8,7$; in both
cases the $T$-diagonals are $0,7,7$ and $7,7,7$ respectively,
all in $S$.) However both candidates fail $B$-closure at the
off-diagonal cell $B(7,8)=9$: in $\{0,7,8\}$, $9\notin S$; in
$\{6,7,8\}$, $9\notin S$. Hence no size-$3$ subset is jointly
closed.

\emph{Sizes $4$ through $10$ --- uniqueness within size class.}
We must confirm that the listed $S_{k}$ is the \emph{unique}
jointly-closed subset of its size class, for $k\in\{4,5,6,7,8,9\}$.
A useful structural reduction: \emph{any jointly-closed subset
$S \subseteq \Zten$ of size at least $4$ contains
$\Cfour = \{0, 7, 8, 9\}$.}

To establish $\Cfour \subseteq S$ for any jointly-closed
$|S| \geq 4$: assume $S \neq \{0\}$ (the size-$1$ case),
and that $S$ contains some element $j_{0} \neq 0$.

\emph{Case 1: $j_{0} \in \{1,\dots,6\}$.} The diagonal
$B(j_{0},j_{0}) = j_{0} + 1 \pmod{10}$ must lie in $S$. Iterating
$j \mapsto B(j,j) = j + 1$ on the consecutive integers
$\{j_{0}, j_{0}+1, \dots, 6, 7\}$ then forces all of these into
$S$ (the diagonal increment chain stops at $j = 7$ since
$B(7,7) = 8$ and $B(8,8) = 7$ stays in $\{7, 8\}$). Hence
$\{7, 8\} \subseteq S$. Now, with $7, 8 \in S$, joint closure
forces $B(7, 8) = 9 \in S$, then $B(7, 9) = 0 \in S$. Therefore
$\Cfour = \{0, 7, 8, 9\} \subseteq S$.

\emph{Case 2: $j_{0} = 7$.} Then $B(7,7) = 8 \in S$, and joint
closure proceeds as in Case~1's tail: $B(7,8) = 9 \in S$,
$B(7,9) = 0 \in S$, giving $\Cfour \subseteq S$.

\emph{Case 3: $j_{0} \in \{8, 9\}$.} Then $B(j_{0}, j_{0}) \in
\{7, 0\}$ lies in $S$. If $7 \in S$, Case~2 applies. If
$j_{0} = 9$ and $0 \in S$ but $7 \notin S$: we need a third
element to reach size $\geq 4$. Off-diagonal closure of $B$ on
$\{0, 9\}$ gives $B(0, 9) = 9$ (already in $S$), so closure on
$\{0, 9\}$ itself is achieved in size $2$, ruled out by the
size-$2$ analysis. Adding any $k \notin \{0, 9\}$ to reach size
$\geq 3$ then triggers Case~1 or Case~2 on $k$. Symmetric
argument from $j_{0} = 8$.

In every case $\Cfour \subseteq S$. Hence the listed $S_{k}$ for
$k \geq 4$ are obtained by adjoining operators
$\{2, 3, 4, 5, 6\}$ to $\Cfour$ in the order forced by joint
closure (cf.\ Remark~\ref{rem:sigma-walk}).

Direct enumeration over the $\binom{10}{k}$ candidates at each
size (totalling $210+252+210+120+45+10=847$ subsets across sizes
$4$–$9$, plus $\binom{10}{10}=1$ at size $10$) confirms that each
non-listed candidate fails $T$- or $B$-closure on at least one
cell, completing the uniqueness verification. The verification
script (\S\ref{sec:verify}) records the failing cell for each
non-jointly-closed candidate.

The full enumeration of all $1023$ non-empty subsets of $\Zten$,
with the distinguishing failing cell recorded for each
non-jointly-closed subset, is performed by the verification
script \texttt{4core\_verification.py} (\S\ref{sec:verify}).
\end{proof}

\begin{remark}[Structural reading of the chain]\label{rem:sigma-walk}
The chain has a clean structural description in terms of a specific
permutation $\pi:\Zten\to\Zten$ defined by
\begin{equation}\label{eq:pi}
\pi=(0)(3)(8)(9)(1\;7\;6\;5\;4\;2)
\quad\text{in cycle notation.}
\end{equation}
The operators $\{0,3,8,9\}$ are $\pi$-fixed; the operators
$\{1,2,4,5,6,7\}$ form a single $6$-cycle. We do not derive
$\pi$ from the tables $T$ and $B$; rather, $\pi$ is introduced
here as a candidate organizing permutation of the operator
labels. (We use the symbol $\pi$ in this paper to avoid a
collision with the symbol $\sigma$ used for the
non-associativity-rate function in the companion
paper~\cite{SandersGish2026Sigma}, which is a different object.
The two are not yet structurally connected: whether the
permutation $\pi$ here arises from the same operator-substrate
construction that defines the $\sigma$-rate function in the
companion is an open question of the substrate framework, and
not a result of the present paper.) The structural reading below
is the observation that this specific $\pi$ also organizes the
chain of Theorem~\ref{thm:chain}. The chain proceeds:
\begin{itemize}\setlength\itemsep{0pt}
\item Step $1\to 4$: adds $\{7,8,9\}$. Of these, $7$ is in the
$6$-cycle and $\{8,9\}$ are $\pi$-fixed.
\item Steps $4\to 5$, $5\to 6$, $6\to 7$: each adds the next
element of the $\pi$-forward orbit of $7$, namely
$\pi(7)=6$, $\pi^{2}(7)=5$, $\pi^{3}(7)=4$.
\item Step $7\to 8$: adds $3$, the remaining $\pi$-fixed
operator. The $\pi$-orbit walk pauses for one step:
$\pi^{4}(7)=2$ would be the next $\pi$-orbit element, but
adding $2$ alone forces $3=B(2,2)$ into the closure as well,
producing a size jump of $2$. Adding $3$ alone closes more
cheaply.
\item Steps $8\to 9$, $9\to 10$: completes the $\pi$-orbit
walk, adding $\pi^{4}(7)=2$ at size $9$ and
$\pi^{5}(7)=1$ at size $10$.
\end{itemize}
The $\pi$-fixed operators $\{0,3,8,9\}$ enter at three
positions: $0$ at the bottom (size $1$), $\{8,9\}$ in the
size-$1\to 4$ jump, and $3$ at the size-$7\to 8$ bridge. The
$\pi$-cycle $\{1,7,6,5,4,2\}$ enters in $\pi$-forward order,
with the size-$8$ bridge interrupting between
$4=\pi^{3}(7)$ and $2=\pi^{4}(7)$.
\end{remark}

\begin{corollary}[Forbidden sizes]\label{cor:forbidden}
No subset of $\Zten$ of size $2$ or size $3$ is jointly closed under
$T$ and $B$.
\end{corollary}

\begin{proof}
Part~(b) of the proof of Theorem~\ref{thm:chain}.
\end{proof}

%% -------- section 4: 4-core --------
\section{The 4-core is jointly closed}\label{sec:4core}

The minimal non-trivial element of the chain in
Theorem~\ref{thm:chain} is the four-element subset
$\Cfour=\{0,7,8,9\}$.

\begin{corollary}[4-core closure]\label{cor:4core}
$\Cfour$ is jointly closed under $T$ and $B$. Explicitly,
$T(\Cfour\times\Cfour)\subseteq\{0,7\}\subset\Cfour$ and
$B(\Cfour\times\Cfour)=\Cfour$.
\end{corollary}

\begin{proof}
Part~(a) of the proof of Theorem~\ref{thm:chain}, displayed
explicitly via the $4\times 4$ restrictions.
\end{proof}

\begin{remark}\label{rem:rank-disparity}
The image of $T\big|_{\Cfour\times\Cfour}$ is the proper subset
$\{0,7\}\subset\Cfour$. The image of $B\big|_{\Cfour\times\Cfour}$
is the full $\Cfour$. The two operations therefore differ
markedly in rank on $\Cfour$. The two restrictions agree on
exactly $4$ of the $16$ cells: $(0,0)\to 0$, $(0,7)\to 7$,
$(7,0)\to 7$, and $(8,8)\to 7$. They disagree on the remaining
$12$ cells. The per-coordinate fuse data of
Proposition~\ref{prop:fuse-data} catalogs the resulting
distribution explicitly.
\end{remark}

%% -------- section 5: per-coordinate fuse data on the 4-core --------
\section{Per-coordinate fuse data on the 4-core}\label{sec:norm}

For $p\in\R^{10}$ supported on $\Cfour$, write
$(v,h,\beta,r)=(p_{0},p_{7},p_{8},p_{9})$.

\begin{proposition}[$\Cfour$-fuse data]\label{prop:fuse-data}
Let $p\in\R^{10}$ be supported on $\Cfour$. Then
$\Tfuse(p)$ and $\Bfuse(p)$ are also supported on $\Cfour$
(consequence of Theorem~\ref{thm:chain}), and their four
$\Cfour$-coordinates are explicitly
\begin{align*}
\Tfuse(p)_{0} &= v^{2}+2v\beta+2vr,
   & \Tfuse(p)_{7} &= 2vh+h^{2}+2h\beta+2hr+\beta^{2}+2\beta r+r^{2},\\
\Tfuse(p)_{8} &= 0,
   & \Tfuse(p)_{9} &= 0,
\end{align*}
\begin{align*}
\Bfuse(p)_{0} &= v^{2}+2hr+r^{2},
   & \Bfuse(p)_{7} &= 2vh+\beta^{2},\\
\Bfuse(p)_{8} &= 2v\beta+h^{2}+2\beta r,
   & \Bfuse(p)_{9} &= 2vr+2h\beta.
\end{align*}
The image of $T\big|_{\Cfour\times\Cfour}$ is $\{0,7\}$
($\Tfuse$ has rank $2$ on $\Cfour$); the image of
$B\big|_{\Cfour\times\Cfour}$ is the full $\Cfour$ ($\Bfuse$ has
rank $4$ on $\Cfour$). Both fuses sum to the same total
$Z_{T}(p) = Z_{B}(p) = (v+h+\beta+r)^{2}$, by the total-mass
identity (Remark~\ref{rem:total-mass}).
\end{proposition}

\begin{proof}
We compute the four 4-core-supported coordinates of $\Tfuse(p)$
and $\Bfuse(p)$ from Definition~\ref{def:fuse} and the restricted
tables of Corollary~\ref{cor:4core}.

For $T$: cells of $T\big|_{\Cfour\times\Cfour}$ producing $0$ are
$(0,0),(0,8),(8,0),(0,9),(9,0)$ ($5$ cells). Cells producing $7$
are the remaining $11$ cells. Tabulating fuse coordinates in
$(v,h,\beta,r)$:
\begin{align*}
\Tfuse(p)_{0}&=v^{2}+2v\beta+2vr,\\
\Tfuse(p)_{7}&=2vh+h^{2}+2h\beta+2hr+\beta^{2}+2\beta r+r^{2},\\
\Tfuse(p)_{8}&=0,\qquad \Tfuse(p)_{9}=0.
\end{align*}
Summing over $\Cfour$ (and noting $\Tfuse(p)_{c}=0$ for
$c\notin\Cfour$ when $p$ is $\Cfour$-supported, since $T(i,j)\in\Cfour$
for $i,j\in\Cfour$):
\[
Z_{T}(p)\;=\;v^{2}+h^{2}+\beta^{2}+r^{2}+2(vh+v\beta+vr+h\beta+hr+\beta r)\;=\;(v+h+\beta+r)^{2}.
\]

For $B$: cells of $B\big|_{\Cfour\times\Cfour}$ producing $0$ are
$(0,0),(7,9),(9,7),(9,9)$. Cells producing $7$ are
$(0,7),(7,0),(8,8)$. Cells producing $8$ are
$(0,8),(8,0),(7,7),(8,9),(9,8)$. Cells producing $9$ are
$(0,9),(9,0),(7,8),(8,7)$. Tabulating fuse coordinates:
\begin{align*}
\Bfuse(p)_{0}&=v^{2}+2hr+r^{2},\\
\Bfuse(p)_{7}&=2vh+\beta^{2},\\
\Bfuse(p)_{8}&=2v\beta+h^{2}+2\beta r,\\
\Bfuse(p)_{9}&=2vr+2h\beta.
\end{align*}
Summing:
\[
Z_{B}(p)\;=\;v^{2}+h^{2}+\beta^{2}+r^{2}+2(vh+v\beta+vr+h\beta+hr+\beta r)\;=\;(v+h+\beta+r)^{2}.
\]
Both equal $(v+h+\beta+r)^{2}$.
\end{proof}

\begin{corollary}[Normalizer is one on the 4-core simplex]\label{cor:Z=1}
For $p$ in the standard $9$-simplex $\Delta^{9}$ with support in
$\Cfour$, $Z_{T}(p)=Z_{B}(p)=1$.
\end{corollary}

\begin{remark}\label{rem:rank-disparity-detailed}
The two operations agree on only $4$ of the $16$ cells of the
$\Cfour\times\Cfour$ restriction. The remaining $12$ cells
distribute mass into different output buckets:
$\Tfuse$ concentrates all mass on $\{0, 7\}$ (rank-$2$ image),
while $\Bfuse$ spreads mass across all of $\Cfour$
(rank-$4$ image). The per-coordinate values of
Proposition~\ref{prop:fuse-data} therefore differ markedly --
e.g., $\Tfuse(p)_{8} = 0$ identically while $\Bfuse(p)_{8} =
2v\beta + h^{2} + 2\beta r$, a polynomial that vanishes only on
the codimension-$1$ subset $\{h = 0,\ \beta(v + r) = 0\}$ of
$\Delta^{3}_{\Cfour}$. The total-mass identity
$Z_{T}(p) = Z_{B}(p)$ is not the substantive structural
observation; the substantive observation is that the
per-coordinate mass distributions $\Tfuse(p)_{c}$ and
$\Bfuse(p)_{c}$ on the four cells $c\in\Cfour$ are explicit
non-negative quadratics in $(v, h, \beta, r)$, and that they
preserve $\Cfour$-support, so that any convex combination
$\alpha\Tfuse + (1-\alpha)\Bfuse$ for $\alpha\in[0,1]$ is a
self-map of the $4$-core simplex $\Delta^{3}_{\Cfour}$. The
dynamical structure of this family is the subject of
\S\S\ref{sec:dynamics}--\ref{sec:alpha-unique} below.
\end{remark}

%% -------- section 6: dynamics on the simplex --------
\section{The convex-combination iteration on the $4$-core simplex}\label{sec:dynamics}

\subsection{Setup.} Let $\Delta^{3}_{\Cfour} := \{p\in\R^{10} :
p_{i}\ge 0,\ \sum_{i\in\Cfour}p_{i} = 1,\ p_{i} = 0\text{ for }
i\notin\Cfour\}$ denote the standard $3$-simplex of probability
distributions supported on $\Cfour = \{0, 7, 8, 9\}$. For
$\alpha\in[0,1]$, define
\[
  F_{\alpha}: \Delta^{3}_{\Cfour}\to \Delta^{3}_{\Cfour},
  \qquad
  F_{\alpha}(p) := \alpha\,\Tfuse(p) + (1-\alpha)\,\Bfuse(p),
\]
where $\Tfuse$ and $\Bfuse$ are normalized fuses (each is divided
by its $L^{1}$-norm $Z_{T}(p)$, $Z_{B}(p)$ respectively, so that
$\Tfuse(p), \Bfuse(p)\in\Delta^{3}_{\Cfour}$ for any
$p\in\Delta^{3}_{\Cfour}$). The map $F_{\alpha}$ is well-defined
on $\Delta^{3}_{\Cfour}$ by Proposition~\ref{prop:fuse-data} (both
$\Tfuse$ and $\Bfuse$ preserve $\Cfour$-support and produce
non-negative outputs on $\Delta^{3}_{\Cfour}$).

\begin{remark}
Since $Z_{T}(p) = Z_{B}(p) = (\sum p_{i})^{2} = 1$ on
$\Delta^{3}_{\Cfour}$, the normalization is trivial on the
simplex (no rescaling is needed). We will therefore take
$F_{\alpha}(p) = \alpha\Tfuse(p) + (1-\alpha)\Bfuse(p)$ directly
in coordinate calculations below, with $\Tfuse, \Bfuse$ in their
unnormalized polynomial forms of
Proposition~\ref{prop:fuse-data}.
\end{remark}

\subsection{Boundary fixed points.} The boundary
$\partial\Delta^{3}_{\Cfour}$ consists of the four vertices
$\{e_{0}, e_{7}, e_{8}, e_{9}\}$, the six edges joining pairs of
vertices, and the four triangular faces. We classify boundary
fixed points of $F_{\alpha}$ for $\alpha\in[0,1]$.

\begin{lemma}[Boundary fixed points of $F_{\alpha}$]\label{lem:boundary}
For $\alpha\in[0, 1]$:
\begin{enumerate}\setlength\itemsep{0pt}
\item[(a)] $e_{0}$ is a fixed point of $F_{\alpha}$ for every
$\alpha$. For $\alpha\in(0, 1)$, $e_{0}$ is a repellor: the
spectral radius of $DF_{\alpha}(e_{0})$ is strictly greater
than $1$.
\item[(b)] $e_{7}$ is a fixed point of $F_{\alpha}$ if and only if
$\alpha = 1$.
\item[(c)] Neither $e_{8}$ nor $e_{9}$ is a fixed point of
$F_{\alpha}$ for any $\alpha\in[0,1]$.
\item[(d)] No edge or face fixed point exists in
$\partial\Delta^{3}_{\Cfour}$ for any $\alpha\in[0, 1]$ other
than the four vertices listed above.
\end{enumerate}
\end{lemma}

\begin{proof}
\textbf{(a)} At $e_{0} = (1, 0, 0, 0)$, the per-coordinate fuse
data of Proposition~\ref{prop:fuse-data} gives
$\Tfuse(e_{0}) = (v^{2}, 0, 0, 0) = e_{0}$ and
$\Bfuse(e_{0}) = (v^{2}, 0, 0, 0) = e_{0}$, so $F_{\alpha}(e_{0})
= e_{0}$ for every $\alpha\in[0,1]$.
For the spectral-radius computation, write the simplex tangent
space at $e_{0}$ in coordinates $(h, \beta, r)$ (with
$v = 1 - h - \beta - r$). Substituting $v = 1 - h - \beta - r$
into the per-coordinate fuse polynomials and reading off the
linear-in-$(h, \beta, r)$ terms gives
\[
  D\Tfuse(e_{0})\big|_{(h,\beta,r)} = \mathrm{diag}(2, 0, 0),
  \qquad
  D\Bfuse(e_{0})\big|_{(h,\beta,r)} = 2\,I_{3}.
\]
For $\Tfuse$: only $\Tfuse(p)_{7} = 2vh + h^{2} + \cdots$ has a
linear-in-tangent contribution, $\partial_{h}\Tfuse(p)_{7}|_{e_{0}}
= 2$; the $\Tfuse(p)_{0} = v^{2} + \cdots$ component contributes
no first-order term in $(h, \beta, r)$ since
$\partial_{h}v^{2}|_{e_{0}} = -2$ but the corresponding
tangent-direction is the $v$-direction, off-simplex. For
$\Bfuse$: each of
$\Bfuse(p)_{7} = 2vh + \beta^{2}$,
$\Bfuse(p)_{8} = 2v\beta + h^{2} + 2\beta r$,
$\Bfuse(p)_{9} = 2vr + 2 h\beta$
contributes $2$ on the diagonal in its respective coordinate, and
the off-diagonal first-order terms vanish at $e_{0}$. Thus
$D\Bfuse(e_{0}) = 2I_{3}$. Combining,
\[
  DF_{\alpha}(e_{0})
  = \alpha \cdot \mathrm{diag}(2, 0, 0) + (1-\alpha)\cdot 2 I_{3}
  = \mathrm{diag}\bigl(2,\ 2(1-\alpha),\ 2(1-\alpha)\bigr),
\]
with spectral radius $\rho(DF_{\alpha}(e_{0})) = 2 > 1$ for every
$\alpha\in[0, 1]$. Hence $e_{0}$ is a (linear) repellor on the
interior of the simplex for every $\alpha\in(0, 1)$.

\textbf{(b)} At $e_{7}$: $T(7, 7) = 7$ gives $\Tfuse(e_{7}) =
e_{7}$, while $B(7, 7) = 8$ gives $\Bfuse(e_{7}) = e_{8}$.
Therefore $F_{\alpha}(e_{7}) = \alpha e_{7} + (1-\alpha) e_{8}$,
fixed iff $\alpha = 1$.

\textbf{(c)} At $e_{8}$: $T(8, 8) = 7$ gives $\Tfuse(e_{8}) =
e_{7}$; $B(8, 8) = 7$ gives $\Bfuse(e_{8}) = e_{7}$. Hence
$F_{\alpha}(e_{8}) = e_{7}\neq e_{8}$. At $e_{9}$:
$T(9, 9) = 7$, $B(9, 9) = 0$; so $F_{\alpha}(e_{9}) =
\alpha e_{7} + (1-\alpha)e_{0}\neq e_{9}$.

\textbf{(d)} We rule out edge and face fixed points by checking
each codimension-$1$ face of $\Delta^{3}_{\Cfour}$.
\begin{itemize}\setlength\itemsep{0pt}
\item Face $\{v = 0\}$, i.e., $h + \beta + r = 1$. From
Proposition~\ref{prop:fuse-data}, $\Tfuse(p)_{0} = 0$ and
$\Bfuse(p)_{0} = r^{2} + 2hr$, so $F_{\alpha}(p)_{0} =
(1-\alpha)(r^{2} + 2hr)$. Fixed-point condition forces
$(1-\alpha)(r^{2} + 2hr) = 0$, i.e., $r(r + 2h) = 0$ for
$\alpha\in[0, 1)$. Since $r, h\geq 0$, this gives $r = 0$.
Reducing to $h + \beta = 1$ (the edge $\{v = r = 0\}$,
analyzed below).
\item Face $\{h = 0\}$. $\Tfuse(p)_{7} = \beta^{2} + 2\beta r +
r^{2}$ and $\Bfuse(p)_{7} = \beta^{2}$, so $F_{\alpha}(p)_{7} =
\alpha(\beta^{2} + 2\beta r + r^{2}) + (1-\alpha)\beta^{2}$.
Fixed-point condition $h = F_{\alpha}(p)_{7}$ with $h = 0$
gives $\alpha(\beta^{2} + 2\beta r + r^{2}) +
(1-\alpha)\beta^{2} = 0$. All three terms are non-negative for
$\alpha\in[0, 1]$; the equation forces $\beta = 0$ and
$\alpha r^{2} = 0$ (so $r = 0$ if $\alpha > 0$). Reducing to
$\{h = \beta = r = 0\}$, the vertex $e_{0}$.
\item Face $\{\beta = 0\}$. $\Tfuse(p)_{8} = 0$ identically, so
$F_{\alpha}(p)_{8} = (1-\alpha)\Bfuse(p)_{8} =
(1-\alpha)(2v\beta + h^{2} + 2\beta r) =
(1-\alpha)h^{2}$ on $\{\beta = 0\}$. Fixed-point condition
$\beta = (1-\alpha)h^{2}$ with $\beta = 0$ forces $h = 0$ for
$\alpha\in[0, 1)$, reducing to $\{h = \beta = 0\}$ on which
$F_{\alpha}(p)_{9} = (1-\alpha)(2vr + 2h\beta) =
2(1-\alpha)vr$ and $r = 2(1-\alpha)vr$. With $v = 1 - r$ this
gives $r = 0$ or $r = (2(1-\alpha) - 1)/(2(1-\alpha))$; the
second value is $< 0$ for $\alpha > 1/2$ and $\geq 1$ for
$\alpha \leq 0$, so ruled out for $\alpha\in[0, 1]$. Reducing
to $e_{0}$.
\item Face $\{r = 0\}$. Symmetric to $\{\beta = 0\}$ by exchange
of $\beta\leftrightarrow r$: $F_{\alpha}(p)_{9} = 2(1-\alpha)
v r + 2(1-\alpha) h\beta$, vanishing on $\{r = 0\}$ iff
$h\beta = 0$, reducing to one of the analyzed faces.
\end{itemize}
The codimension-$2$ edge cases reduce to vertices by the same
arguments. Hence no boundary fixed point exists in
$\partial\Delta^{3}_{\Cfour}$ for any $\alpha\in[0, 1]$ other
than $e_{0}$ (always) and $e_{7}$ (at $\alpha = 1$).
\end{proof}

\subsection{Existence at every $\alpha$; uniqueness at
$\alpha = \tfrac12$.} We separate the easy global existence
statement from the uniqueness statement which we will prove only
at the symmetric mixing weight (the only case used in
\S\S\ref{sec:hbr}--\ref{sec:alpha-unique}).

\begin{theorem}[Interior existence]\label{thm:attractor}
For each $\alpha\in(0,1)$ there exists at least one interior
fixed point $p^*(\alpha)\in\Delta^{3}_{\Cfour}^{\circ}$ of
$F_{\alpha}$.
\end{theorem}

\begin{proof}
The map $F_{\alpha}: \Delta^{3}_{\Cfour}\to\Delta^{3}_{\Cfour}$
is continuous on a compact convex set, so by Brouwer's fixed-point
theorem at least one fixed point exists. By Lemma~\ref{lem:boundary},
no boundary fixed point exists for $\alpha\in(0, 1)$ except
$e_{0}$, which is a repellor; standard topological arguments
(e.g., Lefschetz fixed-point theory or the index calculation for
a repellor on the boundary of a contractible set) then place at
least one fixed point in the interior.
\end{proof}

\begin{remark}[On global uniqueness]\label{rem:uniqueness}
We do not prove that the interior fixed point of
Theorem~\ref{thm:attractor} is unique for every
$\alpha\in(0, 1)$, nor that the iteration globally attracts to
it. Numerical verification
(\texttt{04\_bridge\_attractor.py}, \S\ref{sec:verify}) confirms
both properties at $19$ sampled values
$\alpha\in[0.05, 0.95]$: at each, the spectral radius of
$DF_{\alpha}(p^*(\alpha))$ is strictly less than $1$ and
iteration from random Dirichlet seeds converges to a single
limit. A structural proof of global uniqueness and global
attraction (for instance, via a contraction argument in the
Hilbert (Birkhoff) projective metric on
$\Delta^{3}_{\Cfour}^{\circ}$~\cite{Birkhoff1957, Bushell1973})
is open. The
results of \S\S\ref{sec:hbr}--\ref{sec:alpha-unique} require only
existence (Theorem~\ref{thm:attractor}) and uniqueness at
$\alpha = \tfrac12$, which we prove in
Theorem~\ref{thm:attractor-half} below.
\end{remark}

\begin{theorem}[Uniqueness at $\alpha = \tfrac12$]\label{thm:attractor-half}
The map $F_{1/2}$ has a unique interior fixed point
$p^*\in\Delta^{3}_{\Cfour}^{\circ}$.
\end{theorem}

\begin{proof}
The fixed-point system $F_{1/2}(p) = p$ written in
$(v, h, \beta, r)$ coordinates is a system of three independent
quadratic equations (one each for the $h$, $\beta$, $r$
components, with $v = 1 - h - \beta - r$). Quadratic systems in
three variables admit only finitely many complex solutions
(generically at most $2^{3} = 8$ by B\'ezout's bound). The
explicit closed forms of Theorems~\ref{thm:hbr-ratio}
and~\ref{thm:galois} below pin one such solution exactly: the
positive real solution with $v, h, \beta, r > 0$ summing to $1$,
$h/\beta = 1+\sqrt{3}$, and $r/\beta$ the unique positive real
root of the irreducible quartic $f(x)$. Direct numerical
enumeration of the remaining B\'ezout solutions
(\texttt{04\_bridge\_attractor.py}, \S\ref{sec:verify}) confirms
that no other lies in $\Delta^{3}_{\Cfour}^{\circ}$:
each remaining solution has either a negative coordinate or a
coordinate $> 1$ or fails the closure constraint.
Combined with Theorem~\ref{thm:attractor}, this establishes
that $p^*$ is the unique interior fixed point of $F_{1/2}$.
\end{proof}

%% -------- section 7: closed-form ratio at alpha=1/2 --------
\section{Closed-form ratio $h/\beta = 1+\sqrt{3}$ at $\alpha = \tfrac12$}\label{sec:hbr}

We now restrict to $\alpha = \tfrac12$. Throughout this section
$p^*$ denotes the unique interior fixed point of $F_{1/2}$ on
$\Delta^{3}_{\Cfour}^{\circ}$ from Theorem~\ref{thm:attractor-half}.
Write its coordinates as
$(v, h, \beta, r) = (p^*_{0}, p^*_{7}, p^*_{8}, p^*_{9})$.

\begin{theorem}[$h/\beta$ closed form]\label{thm:hbr-ratio}
At $\alpha = \tfrac12$, the attractor coordinates satisfy
\[
  \left(\frac{h}{\beta}\right)^{2} - 2\,\frac{h}{\beta} - 2 \;=\; 0,
  \qquad \frac{h}{\beta} \;=\; 1 + \sqrt{3}.
\]
\end{theorem}

\begin{proof}
The fixed-point equation $F_{1/2}(p^*) = p^*$ unpacks
coordinate-by-coordinate using the per-coordinate fuse data of
Proposition~\ref{prop:fuse-data}. The
$\beta$-coordinate (i.e., the $c = 8$ entry) gives
\[
  \beta \;=\; \tfrac12\bigl(\Tfuse(p^*)_{8} + \Bfuse(p^*)_{8}\bigr)
        \;=\; \tfrac12\bigl(0 + (2v\beta + h^{2} + 2\beta r)\bigr)
        \;=\; v\beta + \tfrac12 h^{2} + \beta r,
\]
since $\Tfuse(p^*)_{8} = 0$ identically. Rearranging,
\[
  2\beta \;=\; 2v\beta + h^{2} + 2\beta r,
  \qquad\text{i.e.,}\qquad
  2\beta(1 - v - r) \;=\; h^{2}.
\]
Normalization $v + h + \beta + r = 1$ gives $1 - v - r = h + \beta$, so
\[
  2\beta(h + \beta) \;=\; h^{2},
  \qquad
  2\beta^{2} + 2 h\beta - h^{2} \;=\; 0.
\]
Dividing by $\beta^{2}$ yields the quadratic
$(h/\beta)^{2} - 2(h/\beta) - 2 = 0$, whose positive root is
$h/\beta = 1 + \sqrt{3}$.
\end{proof}

\begin{remark}[Numerical check]\label{rem:hbr-numeric}
Iterating $F_{1/2}$ from the uniform initial condition
$p_{0} = (1,0,\dots,0,1,1,1)/4$ on $\Cfour$ converges within
$60$ iterations (double precision) to
$h^* = 0.540195948486216$, $\beta^* = 0.197725440167385$, with
$h^*/\beta^* = 2.732050807568878$ and $1+\sqrt{3} =
2.732050807568877$ (residual $4.4\times 10^{-16}$). The
verification script \texttt{06\_attractor\_closed\_form.py}
(\S\ref{sec:verify}) runs the iteration symbolically in
\texttt{sympy} and confirms the closed form to all algebraic
orders.
\end{remark}

%% -------- section 8: quartic + Galois D_4 --------
\section{The $\xi^*$-quartic and Galois group $D_{4}$}\label{sec:galois}

The remaining algebraic content of the attractor is captured by
the ratio $\xi^* := r/\beta$.

\begin{theorem}[$\xi^*$-quartic]\label{thm:galois}
Let $p^*$ be the unique interior fixed point of $F_{1/2}$ on
$\Delta^{3}_{\Cfour}^{\circ}$ (Theorem~\ref{thm:attractor-half}).
Then $\xi^* = p^*_{9}/p^*_{8}$ is
the unique positive real root of the irreducible monic integer
quartic
\[
  f(x) \;=\; x^{4} + 4x^{3} - x^{2} + 2x - 2.
\]
The Galois group of $f$ over $\Q$ is the dihedral group
$D_{4}$ of order $8$. The number field $K = \Q[x]/(f)$ has
field discriminant $d_{K} = -2^{4}\cdot 3^{2}\cdot 71$,
class number $h_{K} = 1$, signature $(2, 1)$, and is the field
catalogued as LMFDB~$4.2.10224.1$. The Galois closure
$L\supset K$ has degree $8$ over $\Q$, signature $(0, 4)$, class
number $14$, discriminant $2^{8}\cdot 3^{4}\cdot 71^{4}$, and is
LMFDB~$8.0.526936617216.1$. The intermediate field
$\Q(\sqrt{3})\subset L$ is a genuine subfield of $L$, with
explicit factorization
\[
  f(x) \;=\; \bigl(x^{2}+(2-\sqrt{3})x+(\sqrt{3}-1)\bigr)
             \bigl(x^{2}+(2+\sqrt{3})x-(\sqrt{3}+1)\bigr)
\]
in $\Q(\sqrt{3})[x]$.
\end{theorem}

\begin{proof}
\textit{Derivation through $\Q(\sqrt{3})$.} The substantive
algebraic step is short. The fixed-point equation for the
$r$-coordinate at $\alpha = \tfrac12$ gives, using
Proposition~\ref{prop:fuse-data} and the vanishing
$\Tfuse(p^*)_{9} = 0$,
\[
  r \;=\; \tfrac12\bigl(\Tfuse(p^*)_{9} + \Bfuse(p^*)_{9}\bigr)
        \;=\; \tfrac12(2vr + 2 h\beta)
        \;=\; vr + h\beta.
\]
Substituting $v = 1 - h - \beta - r$ and rearranging,
\[
  r^{2} + (h + \beta)\, r \;-\; h\beta \;=\; 0.
\]
Theorem~\ref{thm:hbr-ratio} gives $h = (1+\sqrt{3})\beta$. Setting
$\xi = r/\beta$ and dividing by $\beta^{2}$,
\[
  \xi^{2} + (2+\sqrt{3})\,\xi \;-\; (1+\sqrt{3}) \;=\; 0
  \qquad\text{over }\Q(\sqrt{3}).
\]
The positive root $\xi^* > 0$ is therefore an algebraic number of
degree $2$ over $\Q(\sqrt{3})$.

\textit{Norm to $\Q$.} The minimal polynomial of $\xi^*$ over $\Q$
is the product of the displayed quadratic with its
$\mathrm{Gal}(\Q(\sqrt{3})/\Q)$-conjugate
(replacing $\sqrt{3}$ with $-\sqrt{3}$):
\begin{align*}
  f(x)
  &= \bigl(x^{2} + (2+\sqrt{3})x - (1+\sqrt{3})\bigr)
     \bigl(x^{2} + (2-\sqrt{3})x - (1-\sqrt{3})\bigr) \\
  &= x^{4} + 4x^{3} - x^{2} + 2x - 2.
\end{align*}
The factorization above is the
$\Q(\sqrt{3})$-factorization of $f$, exhibited as part of the
derivation rather than verified post-hoc. The product equals
the displayed quartic by direct expansion (verifiable in any
computer algebra system; \texttt{07\_full\_closed\_form.py},
\S\ref{sec:verify}).

\textit{Irreducibility over $\Q$.} By Gauss's lemma, any
factorization of $f$ in $\Q[x]$ gives a factorization in
$\Z[x]$. Direct evaluation rules out integer roots:
$f(\pm 1) \in \{4, -8\}$ and $f(\pm 2) \in \{46, -26\}$, none
zero. For an integer factorization
$f = (x^{2} + ax + b)(x^{2} + cx + d)$, matching coefficients
gives the system
\[
  a + c = 4,\quad ac + b + d = -1,\quad
  ad + bc = 2,\quad bd = -2.
\]
The constraint $bd = -2$ over $\Z$ has solutions
$(b, d) \in \{(1, -2), (-1, 2), (2, -1), (-2, 1)\}$. In each
case the linear equation $ad + bc = 2$ together with $c = 4 - a$
determines $a$, and the quadratic $ac + b + d = -1$ gives a
consistency check:
\begin{itemize}\setlength\itemsep{0pt}
\item $(b, d) = (1, -2)$: $-2a + (4 - a) = 2$ gives $a = 2/3
\notin \Z$.
\item $(b, d) = (-1, 2)$: $2a - (4 - a) = 2$ gives $a = 2 \in \Z$,
$c = 2$, but then $ac = 4 \neq -2 = -1 - b - d$, contradicting
the quadratic constraint.
\item $(b, d) = (2, -1)$: $-a + 2(4 - a) = 2$ gives $a = 2 \in \Z$,
$c = 2$, again $ac = 4 \neq -2$.
\item $(b, d) = (-2, 1)$: $a - 2(4 - a) = 2$ gives $a = 10/3
\notin \Z$.
\end{itemize}
No integer factorization exists; $f$ is irreducible over $\Z$,
hence over $\Q$. (For an alternative one-line proof via
reduction, $f$ is irreducible mod $7$; note however that $f$ is
\emph{reducible} mod $5$ as
$f \equiv (x^{2} - 2)(x^{2} - x + 1) \pmod{5}$, so a reduction
argument requires care in the choice of prime.)

\textit{Galois group.} The polynomial discriminant
$\Delta_{f} = -2^{6}\cdot 3^{2}\cdot 71$ (computed directly).
The resolvent cubic of $f$ is $g(y) = y^{3} + y^{2} + 16y + 36
= (y + 2)(y^{2} - y + 18)$, with the unique rational root
$y = -2$ and an irreducible quadratic factor of discriminant
$-71$. By the standard quartic-Galois-group classification (Cohen,
\textit{A Course in Computational Algebraic Number Theory},
\S6.3.2~\cite{Cohen1993}), the existence of exactly one rational
root in the resolvent cubic places
$\mathrm{Gal}(f/\Q)\in\{C_{4}, D_{4}\}$. The two are
distinguished by whether $f$ remains irreducible over
$\Q(\sqrt{\Delta_{f}}) = \Q(\sqrt{-71})$: irreducibility there
gives $D_{4}$, factorization gives $C_{4}$. Direct computation
(sympy \texttt{factor(f, extension=[sqrt(-71)])}) returns $f$
itself as a single irreducible factor, confirming
$\mathrm{Gal}(f/\Q) = D_{4}$.

The $D_{4}$-Galois group is also visible directly from the
derivation above. The closed form $h/\beta = 1+\sqrt{3}$
(Theorem~\ref{thm:hbr-ratio}) generates the quadratic subfield
$\Q(\sqrt{3})\subset\Q(\xi^*)$, and the second nontrivial
quadratic subfield $\Q(\sqrt{-71})$ enters through the
discriminant; the splitting field is the compositum
$\Q(\sqrt{3}, \sqrt{-71}, \xi^*)$, a $D_{4}$-extension of $\Q$.

\textit{Number-field identification.} The field $K = \Q[x]/(f)$
has degree $4$, signature $(2, 1)$ (counting real roots of $f$:
two real, one complex conjugate pair).
The Tschirnhaus substitution $x \mapsto -x - 1$ converts $f$ to
$h(x) = x^{4} - 7x^{2} - 12x - 8$, the canonical defining
polynomial recorded in LMFDB's
catalogue~\cite{LMFDB-4-2-10224-1} for the
quartic field with $d_{K} = -10224 = -2^{4}\cdot 3^{2}\cdot 71$,
class number $h_{K} = 1$, regulator
$R_{K} \approx 8.617$, and Galois group $D_{4}$. Hence
$K = $ LMFDB~$4.2.10224.1$. The Galois closure is the
$D_{4}$-extension recorded
in~\cite{LMFDB-8-0-526936617216-1};
intermediate subfields include $\Q(\sqrt{3})$, $\Q(\sqrt{-71})$,
and the biquadratic $\Q(\sqrt{3}, \sqrt{-71})$.
\end{proof}

\begin{remark}[Novelty]\label{rem:novelty}
The number field $K = $ LMFDB~$4.2.10224.1$ is catalogued and
not new; what is novel is the derivation. The polynomial form
$f(x) = x^{4} + 4x^{3} - x^{2} + 2x - 2$ does not appear in
OEIS for its coefficient sequence $[1, 4, -1, 2, -2]$ as a
distinguished sequence, and we have not located it as a defining
polynomial in a published context tied to a finite-magma
fuse-iteration. The pitch is therefore: a new route to a known
number field, with the ring of definition realized as the
fixed-point coordinates of the symmetric-mixing-weight iteration
$F_{1/2}$ on the $4$-core simplex $\Delta^{3}_{\Cfour}$.
\end{remark}

%% -------- section 9: alpha-uniqueness PSLQ --------
\section{$\alpha$-uniqueness on the Stern--Brocot grid}\label{sec:alpha-unique}

Theorems~\ref{thm:hbr-ratio} and~\ref{thm:galois} exhibit the
attractor $p^*(\tfrac12)$ as a point with both ratios
$h/\beta$ and $r/\beta$ algebraic over $\Q$ of small degree and
small coefficients. We ask: is $\alpha = \tfrac12$ the unique
rational mixing weight with this property?

\subsection{Stern--Brocot scan.}
Let $\mathcal{G}_{\le 7} = \{p/q\in(0,1)\cap\Q : \gcd(p, q) = 1,
\, q\le 7\}$ denote the $17$-point Stern--Brocot grid of
rationals with denominator at most $7$:
\[
  \mathcal{G}_{\le 7} = \bigl\{
    \tfrac{1}{7}, \tfrac{1}{6}, \tfrac{1}{5}, \tfrac{1}{4},
    \tfrac{2}{7}, \tfrac{1}{3}, \tfrac{2}{5}, \tfrac{3}{7},
    \tfrac{1}{2}, \tfrac{4}{7}, \tfrac{3}{5}, \tfrac{2}{3},
    \tfrac{5}{7}, \tfrac{3}{4}, \tfrac{4}{5}, \tfrac{5}{6},
    \tfrac{6}{7}
  \bigr\}.
\]
For each $\alpha\in\mathcal{G}_{\le 7}$ we compute
$p^*(\alpha)$ to $50$-digit precision (\texttt{mpmath}) and apply
the PSLQ algorithm of Ferguson--Bailey~\cite{FergusonBailey1999}
to the bases
$\{1, x, x^{2}, \dots, x^{d}\}$ for $d\in\{2, \dots, 8\}$, with
coefficient bound $|c_{i}|\le 50$ and tolerance
$10^{-42}$. The script
\texttt{alpha\_pslq\_sweep.py} (\S\ref{sec:verify}) reports the
lowest-degree relation found (or no relation, if PSLQ does not
detect any).

\subsection{Result.}
At $\alpha = \tfrac12$, PSLQ recovers exactly:
\begin{align*}
  h/\beta &: \quad x^{2} - 2x - 2 = 0 \quad
    (\text{degree }2,\ \sup|c_{i}| = 2),\\
  r/\beta &: \quad x^{4} + 4x^{3} - x^{2} + 2x - 2 = 0 \quad
    (\text{degree }4,\ \sup|c_{i}| = 4),
\end{align*}
matching Theorems~\ref{thm:hbr-ratio} and~\ref{thm:galois}
exactly. At each of the other $16$ rationals
$\alpha\in\mathcal{G}_{\le 7}\setminus\{\tfrac12\}$, PSLQ finds
\emph{no} integer-coefficient polynomial relation at degree
$\le 8$ with $|c_{i}| \le 50$ in either ratio.

\begin{remark}[Algebraicity is automatic; degree is the question]
\label{rem:degree-bounded}
For any $\alpha\in\Q\cap(0, 1)$, every interior fixed point of
$F_{\alpha}$ has coordinates algebraic over $\Q$, since the
fixed-point system is polynomial with rational coefficients
(after clearing the rational $\alpha$). The interesting empirical
question is therefore not algebraicity but \emph{minimal-polynomial
degree}: at $\alpha = \tfrac12$ the ratios $h/\beta$ and $r/\beta$
have minimal polynomials of degree $2$ and $4$ respectively over
$\Q$, with small integer coefficients, while at the other $16$
rationals in $\mathcal{G}_{\le 7}$ no relation of degree $\le 8$
with coefficients $\le 50$ is detected by PSLQ at $50$-digit
precision.
\end{remark}

\begin{conjecture}[$\alpha$-uniqueness]\label{conj:alpha-unique}
The mixing weight $\alpha = \tfrac12$ is the unique
$\alpha\in\Q\cap(0, 1)$ at which both ratios
$h(\alpha)/\beta(\alpha)$ and $r(\alpha)/\beta(\alpha)$ of an
interior fixed point of $F_{\alpha}$ have minimal polynomials
over $\Q$ of degree at most $4$ and integer coefficients of
absolute value at most $50$.
\end{conjecture}

The above scan over $\mathcal{G}_{\le 7}$ is the empirical
evidence; a structural proof remains open. The natural strategy
would be to derive, for general rational $\alpha = p/q$, the
fixed-point polynomial system in $(v, h, \beta, r, p, q)$ and
characterize specializations $(p, q)$ for which the resulting
algebraic numbers have low minimal-polynomial degree. We have
not extracted a closed-form characterization.

\begin{remark}[Privileged status of $\alpha = \tfrac12$]
At $\alpha = \tfrac12$ the iteration $F_{\alpha}$ is invariant
under the relabeling $(T, B; \alpha)\mapsto (B, T; 1-\alpha)$,
which trivially swaps the two summands and fixes the symmetric
mixing weight. This is a relabeling rather than a structural
automorphism of the data: $T$ and $B$ differ in rank, image,
and left-margin behavior (\S\ref{sec:setup}), and the attractor
$p^*$ at $\tfrac12$ does not in general satisfy a $T \leftrightarrow B$
symmetry beyond the relabeling-fixed condition $\alpha =
1-\alpha$. We do not claim the relabeling explains
the algebraicity-degree drop at $\tfrac12$; a structural
explanation, if one exists, must come from a deeper feature of
the operator-substrate construction yielding $(T, B)$
\cite{SandersGish2026Sigma}, and is open.
\end{remark}

%% -------- section 10: verification --------
\section{Verification}\label{sec:verify}

The results of this paper are independently verified by a bundle
of Python scripts archived in the shared deposit at
DOI~\texttt{10.5281/zenodo.18852047}.

\begin{enumerate}\setlength\itemsep{0pt}
\item \texttt{4core\_verification.py} (\S\S\ref{sec:chain}--\ref{sec:norm}):
joint-closure enumeration over all $1023$ non-empty subsets of
$\Zten$, recording the failing cell for each non-jointly-closed
candidate; symbolic expansion of $Z_{T}(p)$ and $Z_{B}(p)$ on
$\Cfour$ via \texttt{sympy}; symbolic per-coordinate fuse data.

\item \texttt{04\_bridge\_attractor.py} (\S\ref{sec:dynamics}):
direct iteration of $F_{\alpha}$ from uniform and from $50$
random Dirichlet initial conditions on $\Delta^{3}_{\Cfour}$,
confirming convergence to a unique fixed point in
$\le 110$ iterations at double precision and reading off
spectral radii of $DF_{\alpha}$ at the fixed point.

\item \texttt{06\_attractor\_closed\_form.py} (\S\ref{sec:hbr}):
symbolic derivation of the $h/\beta$ quadratic
$(h/\beta)^{2} - 2(h/\beta) - 2 = 0$ from the per-coordinate
fixed-point system, with closed-form solution
$h/\beta = 1 + \sqrt{3}$ at machine precision.

\item \texttt{07\_full\_closed\_form.py} (\S\ref{sec:galois}):
symbolic derivation of the $\xi$-quartic
$f(x) = x^{4} + 4x^{3} - x^{2} + 2x - 2$; computation of
$\mathrm{disc}(f) = -2^{6}\cdot 3^{2}\cdot 71$ and the resolvent
cubic $g(y) = y^{3} + y^{2} + 16y + 36$ with rational-root
extraction; verification of $D_{4}$ Galois group by checking
irreducibility of $f$ over $\Q(\sqrt{-71})$;
$\Q(\sqrt{3})$-factorization via
\texttt{sympy.factor(f, extension=[sqrt(3)])}; LMFDB cross-check
that the Tschirnhaus substitution $x \mapsto -x - 1$ converts $f$
to LMFDB's defining polynomial $x^{4} - 7x^{2} - 12x - 8$.

\item \texttt{alpha\_pslq\_sweep.py} (\S\ref{sec:alpha-unique}):
$50$-digit \texttt{mpmath} iteration to fixed point at each
$\alpha\in\mathcal{G}_{\le 7}$; PSLQ scan at degrees $2$--$8$
with coefficient bound $50$; recovery of the quadratic and
quartic at $\alpha = \tfrac12$, no detection at the other $16$
rationals.
\end{enumerate}

All scripts run on \texttt{numpy}~+~\texttt{sympy}~+~\texttt{mpmath} alone;
each completes in under one minute on a standard laptop. The
verification output is deterministic.

%% -------- section 7: scope --------
\section{Scope and limits}\label{sec:scope}

The results of this paper are about the specific tables $T$ and
$B$ of Definition~\ref{def:tables}. We do \emph{not} claim:
\begin{itemize}\setlength\itemsep{0pt}
\item that the chain structure of Theorem~\ref{thm:chain} is
shared by other natural pairs of binary operations on $\Zten$;
the chain is a property of the specific $(T, B)$ given.
\item that the total-mass identity
$Z_{T}(p) = Z_{B}(p) = (v + h + \beta + r)^{2}$ is a
non-trivial coincidence: as Remark~\ref{rem:total-mass} states,
$Z_{M}(p) = (\sum_{i}p_{i})^{2}$ holds for any binary operation
$M$ on $\Zten$ by elementary bilinearity. The substantive
structural content of \S\ref{sec:norm} is in the explicit
per-coordinate fuse data of Proposition~\ref{prop:fuse-data}
(rank disparity, symbolic forms) and the support-preservation
property inherited from joint closure.
\item that the closed-form ratio $h/\beta = 1+\sqrt{3}$ at
$\alpha = \tfrac12$ generalizes to other rational mixing
weights or to other natural pairs of operations on $\Zten$. The
PSLQ Stern--Brocot scan over $\mathcal{G}_{\le 7}$ (\S\ref{sec:alpha-unique})
empirically suggests $\alpha = \tfrac12$ is privileged within the
present pair, but a structural proof of
Conjecture~\ref{conj:alpha-unique} is open.
\item that the number field LMFDB~$4.2.10224.1$ is itself new;
it is catalogued. What is novel is the realization of this field
as the ring of definition of the symmetric-mixing-weight
attractor of $F_{1/2}$ on $\Delta^{3}_{\Cfour}$, with the
$\Q(\sqrt{3})$-subfield arithmetically anchoring
$h/\beta = 1+\sqrt{3}$.
\item that the operator-substrate construction yielding $(T, B)$
from a deeper geometric structure is treated here; the present
paper takes $(T, B)$ at $N = 10$ as given and refers to the
companion paper~\cite{SandersGish2026Sigma} for that construction.
\end{itemize}

We \emph{do} claim
Theorems~\ref{thm:chain},~\ref{thm:attractor},~\ref{thm:attractor-half},
\ref{thm:hbr-ratio}, and~\ref{thm:galois},
Proposition~\ref{prop:fuse-data}, and
Conjecture~\ref{conj:alpha-unique} (as a conjecture, with the
empirical PSLQ-grid evidence detailed in
\S\ref{sec:alpha-unique}). All are verified at machine precision
or symbolically by the scripts catalogued in \S\ref{sec:verify};
the scripts run in $\le 1$ minute on standard hardware and
produce deterministic output.

The structural $\pi$-walk reading
(Remark~\ref{rem:sigma-walk}) is offered as commentary on the
chain of Theorem~\ref{thm:chain}; the permutation $\pi$ is
introduced in this paper and is not derived from $(T, B)$
themselves.

%% -------- section 8: forward directions --------
\section{Forward directions}\label{sec:forward}

The pair $(T, B)$ on $\Zten$ admits further structural phenomena
treated in companion work in preparation:

\paragraph{$\F_{p}$-universality.} The chain of
Theorem~\ref{thm:chain} and the per-coordinate fuse data of
Proposition~\ref{prop:fuse-data} persist when $(T, B)$ at
$N = 10$ are replaced by analogous tables $(T_{p}, B_{p})$
constructed via the same operator-substrate recipe over
$\F_{p}$-substrates of comparable size, for
$p \in \{2, 3, 5, 7, 11, 13\}$. The construction of
$(T_{p}, B_{p})$ from $\F_{p}$ and a structural proof of the
$\F_{p}$-universality is the subject of a companion bridge-sprint
paper in
preparation~\cite{SandersClaudeChat2026BridgeSprint}.

\paragraph{Structural proof of Conjecture~\ref{conj:alpha-unique}.}
The natural strategy is to derive, for general rational
$\alpha = p/q$, the polynomial system in
$(v(\alpha), h(\alpha), \beta(\alpha), r(\alpha), p, q)$
governing the attractor coordinates, and to characterize
specializations $(p, q)$ for which the system admits a
low-degree algebraic relation. The Stern--Brocot empirical
evidence (\S\ref{sec:alpha-unique}) is consistent with
$\alpha = \tfrac12$ being the unique rational case; we have not
extracted a closed-form characterization.

\paragraph{Operator-substrate construction.}
The operator-substrate construction yielding $(T, B)$ from a
deeper geometric structure is detailed in the companion
paper~\cite{SandersGish2026Sigma}.

%% -------- references --------
\begin{thebibliography}{99}

\bibitem{DrapalKepka1985}
A.~Dr\'{a}pal and T.~Kepka.
\newblock Group distances of Latin squares.
\newblock \emph{Commentationes Mathematicae Universitatis Carolinae},
26(2):275--283, 1985.

\bibitem{Kepka1980}
T.~Kepka.
\newblock A note on associative triples of elements in cancellation groupoids.
\newblock \emph{Commentationes Mathematicae Universitatis Carolinae},
21(3):479--487, 1980.

\bibitem{DrapalWanless2021}
A.~Dr\'{a}pal and I.~M.~Wanless.
\newblock Maximally non-associative quasigroups.
\newblock \emph{Journal of Combinatorial Theory, Series A},
184:105510, 2021.

\bibitem{Birkhoff1957}
G.~Birkhoff.
\newblock Extensions of {J}entzsch's theorem.
\newblock \emph{Transactions of the American Mathematical Society},
85:219--227, 1957.

\bibitem{Bushell1973}
P.~J.~Bushell.
\newblock Hilbert's metric and positive contraction mappings in a {Banach} space.
\newblock \emph{Archive for Rational Mechanics and Analysis},
52:330--338, 1973.

\bibitem{SandersGish2026Sigma}
B.~R.~Sanders and M.~Gish.
\newblock Non-Associativity Decay in Binary Composition Tables over
$\mathbb{Z}/N\mathbb{Z}$.
\newblock Preprint, 2026.
\newblock Companion paper; archived in shared Zenodo deposit DOI
10.5281/zenodo.18852047.

\bibitem{SandersClaudeChat2026BridgeSprint}
B.~R.~Sanders et al.
\newblock Bridge-Sprint Working Material: $\F_{p}$-Universality of the
Operator-Substrate Construction.
\newblock In preparation, 2026.

\bibitem{Cohen1993}
H.~Cohen.
\newblock \emph{A Course in Computational Algebraic Number Theory},
Graduate Texts in Mathematics 138.
\newblock Springer-Verlag, Berlin, 1993.

\bibitem{FergusonBailey1999}
H.~R.~P.~Ferguson, D.~H.~Bailey, and S.~Arno.
\newblock Analysis of {PSLQ}, an integer relation finding algorithm.
\newblock \emph{Mathematics of Computation}, 68(225):351--369, 1999.

\bibitem{LMFDB-4-2-10224-1}
The {LMFDB} Collaboration.
\newblock Number field 4.2.10224.1.
\newblock The {L}-functions and modular forms database, 2024.
\newblock \url{https://www.lmfdb.org/NumberField/4.2.10224.1}.

\bibitem{LMFDB-8-0-526936617216-1}
The {LMFDB} Collaboration.
\newblock Number field 8.0.526936617216.1.
\newblock The {L}-functions and modular forms database, 2024.
\newblock \url{https://www.lmfdb.org/NumberField/8.0.526936617216.1}.

\end{thebibliography}

\end{document}
